\documentclass[11pt,a4paper]{article}
\usepackage[T1]{fontenc}
\usepackage[utf8]{inputenc}
\usepackage{lmodern,microtype}
\usepackage[margin=25mm,top=24mm,bottom=24mm,headheight=14pt]{geometry}
\usepackage{amsmath,amssymb,amsthm,mathtools,booktabs,array,enumitem,listings,xcolor,fancyhdr,pict2e}
\definecolor{ink}{HTML}{183747}
\definecolor{accent}{HTML}{0A635C}
\definecolor{quiet}{HTML}{4B6269}
\usepackage[hypertexnames=false,colorlinks=true,urlcolor=accent,linkcolor=accent,citecolor=accent,pdfauthor={Shaitetsu},pdftitle={Distance modes and extremal graph inequalities}]{hyperref}
\newtheorem{theorem}{Theorem}[section]
\newtheorem{lemma}[theorem]{Lemma}
\newtheorem{corollary}[theorem]{Corollary}
\newtheorem{proposition}[theorem]{Proposition}
\newcommand{\diam}{\operatorname{diam}}
\newcommand{\rad}{\operatorname{rad}}
\newcommand{\dist}{\operatorname{dist}}
\newcommand{\doi}[1]{\href{https://doi.org/#1}{doi: #1}}
\setlist{topsep=3pt,itemsep=1pt,parsep=0pt}
\lstset{basicstyle=\ttfamily\footnotesize,columns=fullflexible,keepspaces=true,showstringspaces=false,breaklines=true}
\pagestyle{fancy}\fancyhf{}\fancyhead[L]{\small\color{quiet}Distance modes and extremal graph inequalities}
\fancyhead[R]{\small\color{quiet}Research note}\fancyfoot[C]{\thepage}
\renewcommand{\headrulewidth}{0.3pt}\setlength{\emergencystretch}{2em}\raggedbottom
\begin{document}\thispagestyle{plain}
\noindent{\small\sffamily\bfseries\color{accent} GRAPH THEORY\quad /\quad EXTREMAL DISTANCES}\par
\vspace{9pt}\noindent{\LARGE\bfseries\color{ink} Distance modes and extremal graph inequalities}\par
\vspace{5pt}\noindent{\large A sharp correction to Graffiti and exact Wiener extrema at diameter five}\par
\vspace{8pt}\noindent{\small Shaitetsu\quad$\cdot$\quad 21 November 2026}\par
\vspace{9pt}\hrule\vspace{9pt}
\noindent\textbf{Abstract.}
We refute Fajtlowicz's proposed upper bound on the modal distance by the sum of the mean distance and the Randi\'c index. A tree on 157 vertices has unique mode 113 and a sum strictly below 113. For fixed nonnegative coefficients, a universal upper bound by $a$ times the mean distance, $b$ times the Randi\'c index and $c$ times the matching number holds exactly when $b+c\ge2$. We also construct counterexamples of maximum degree three. Independently, we determine the exact maximum Wiener index and all maximizing trees of every order at diameter five by comparing a bounded number of integer profiles. Other results classify the exceptional graphs in two Graffiti radius inequalities, sharpen the harmonic-index gap above mean distance, and characterize two further graph optimization problems.

\medskip\noindent\textbf{Scope.}
Graphs are finite, simple and undirected. Distance assertions concern connected graphs with at least two vertices unless stated otherwise. The Randi\'c index is $R=\sum_{uv\in E}1/\sqrt{d_ud_v}$, the harmonic index is $H=\sum_{uv\in E}2/(d_u+d_v)$, and the mean distance $\mu$ averages over distinct unordered pairs. The mode is the most frequent distance. A change from a mean to a mode changes the truth of the inequality even on trees. We first treat this difference by counting pairs. The Wiener calculation subsequently moves single leaves between branches, and the harmonic-gap estimates use earlier bounds from the optimized quadratic distance energy.

\medskip\noindent\textbf{Guide to the results.}
Sections~\ref{sec:mode-example}--\ref{sec:sharp-mode} give the counterexamples and the sharp replacement.
Section~\ref{sec:diamfive} gives the all-orders Wiener formula.
Sections~\ref{sec:radius-certified}--\ref{sec:odd-fork} resolve the radius and temperature questions.
Sections~\ref{sec:sixth}--\ref{sec:hgap} quantify the gap $H-\mu$ away from stars.
The final two sections concern optimal point sets on block graphs and clique-free subgraphs of threshold graphs.

\section{A tree refuting the modal-distance conjecture}\label{sec:mode-example}
Fajtlowicz's 1988 list of open conjectures, item~5, asks whether the most
frequent distance $m$ in a connected graph always satisfies $m\le\mu+R$
\cite[p.~115]{mode-faj}. Fajtlowicz chooses the largest mode in a tie and
uses distinct vertices for the average \cite[pp.~114--115]{mode-faj}.
This is entry~8 on printed page~14 of
\emph{Written on the Wall} \cite{cat-wow}. The catalogue records a refutation
by Hi Dong Qi in March 1990 for the adjacent entry~9 (item~6 in 1988),
whose right-hand side uses the matching number. The catalogue gives no
edges or construction for Qi's example, so its Randi\'c index cannot
be recovered from that record. Caporossi's thesis proposes a sharpening of entry~8
for $n>3$, namely $\mu+R-m\ge\sqrt{n-1}-2/n$
\cite[p.~33, Conjecture~5, equations~(1.14)--(1.15)]{cat-thesis}.
A tree with a unique mode refutes both assertions.

\begin{theorem}\label{thm:finite-mode}
There exists a tree on 245 vertices with unique modal distance $173$,
mean distance $1070749/14945<72$, and Randi\'c index strictly below~$100$.
\end{theorem}
\begin{proof}
Take a path $v_0,\ldots,v_L$ with $L=170$. Attach vertices $a,b$ to
$v_0$. Attach $t=24$ new leaves to each of $a,b,v_L$. Write $A,B,C$
for the three leaf sets, and $K$ for all remaining vertices, namely the
spine and $a,b$. Thus $n=L+3+3t=245$. The diameter is $L+3=173$.
Precisely the $2t^2=1152$ pairs from $(A\cup B)\times C$ are diametral.

At any other distance, pairs of leaves contribute at most
$\max\{3\binom{t}{2},t^2\}=828$: the three within-set counts occur at
distance two and the $A,B$ count at distance four. For any fixed leaf,
each distance occurs at most twice among its distances to $K$. Indeed
all distances along the spine are distinct; from a near leaf, $a,b$ are
at distances one and three, and from a far leaf their shared distance
is $L+2$, beyond every spine distance. Hence leaf-to-$K$ pairs contribute
at most $6t=144$ per distance. Within $K$, at most $L$ spine pairs occur
at any positive distance, at most two further pairs have one endpoint in
$\{a,b\}$ and one on the spine, and one pair is $\{a,b\}$. Therefore every
other distance has frequency at most $828+144+173=1145<1152$.
The mode is unique. Ordered distinct pairs give the same mode; including
the $245$ diagonal zeros also preserves it.

The degree types of the edges give
\[
 R=\frac{L-2}{2}+\frac{3t}{\sqrt{t+1}}+
   \frac{2}{\sqrt{3(t+1)}}+\frac1{\sqrt6}+
   \frac1{\sqrt{2(t+1)}}
 =84+\frac{72}{5}+\frac{2}{5\sqrt3}+\frac1{\sqrt6}+\frac1{5\sqrt2}<100.
\]
For the strict bound the last three terms are respectively below
$1/4,1/2,1/5$. To compute the distance sum $W$, each tree edge separating
$s$ vertices from $n-s$ contributes $s(n-s)$. The $3t$ pendant edges
have $s=1$, the two edges at the near hubs have $s=t+1$, and the
$j$th spine edge has $s_j=2t+2+j$, starting from the near end. Thus
\[
 W=3t(n-1)+2(t+1)(n-t-1)+\sum_{j=1}^{L}s_j(n-s_j)=2141498.
\]
Dividing by $\binom{245}2=29890$ gives the stated mean. In particular
$R+\mu<172<173=m$. If the average includes diagonal distances, its
value decreases and the same counterexample applies.
The maximum matching has size $88$: deleting $a,b,v_L$ leaves a
170-vertex path and isolated vertices, so a matching uses at most
$85+3$ edges. A maximum matching of this path and one leaf edge at each
of the three deleted hubs attain that number. Thus $\mu+88<160<m$
also verifies the matching-number violation on this tree.
\end{proof}

\subsection{A smaller certificate checked by direct distances}
In the preceding construction use $L=110$, attach 12 leaves to each of the two near hubs and 20 to the remote hub. This gives $n=157$. The schematic labels each group with its multiplicity.

\begin{center}
\setlength{\unitlength}{0.85mm}
\begin{picture}(145,34)
\put(26,17){\line(1,0){63}}\put(26,17){\line(-2,1){18}}\put(26,17){\line(-2,-1){18}}
\put(8,26){\circle*{2}}\put(8,8){\circle*{2}}\put(26,17){\circle*{2}}\put(89,17){\circle*{2}}
\put(8,26){\line(-1,1){4}}\put(8,26){\line(-1,0){5}}\put(8,26){\line(-1,-1){4}}
\put(8,8){\line(-1,1){4}}\put(8,8){\line(-1,0){5}}\put(8,8){\line(-1,-1){4}}
\put(89,17){\line(1,1){5}}\put(89,17){\line(1,0){6}}\put(89,17){\line(1,-1){5}}
\put(11,26){\small 12 leaves}\put(11,4){\small 12 leaves}
\put(37,20){\small path of 110 edges}\put(99,16){\small 20 leaves}
\end{picture}
\end{center}
For a bin of distance $j$ write $A_j$ for the number of unordered pairs. Direct counting as in the first proof gives the five largest bin counts:
\[
 \begin{array}{c|rrrrr}j&113&2&4&3&1\\\hline A_j&480&478&297&178&156\end{array}.
\]
For an independent check, start a breadth-first search at each vertex and count only endpoints of larger index: the histogram totals $12{,}246$ pairs and its weighted sum is $576{,}558$. The same sum follows from $44(156)+2(13)(144)+\sum_{z=21}^{130}z(157-z)$. Hence $\mu=96093/2041<47.09$. Summing the edge weights gives
\[
 R=54+\frac1{\sqrt6}+\frac1{\sqrt{42}}+\frac2{\sqrt{39}}+
       \frac{24}{\sqrt{13}}+\frac{20}{\sqrt{21}}<65.91 .
\]
For a rational check, replace each $\sqrt j$ in these denominators by
$\lfloor 10^4\sqrt j\rfloor/10^4$, evaluated by integer square root. The resulting rational upper bound is smaller than $65.904$. Consequently $\mu+R<113=m$. No minimality of this witness order is needed or asserted.

\section{The sharp coefficient in the distance-mode inequalities}\label{sec:sharp-mode}
Write $m(G)$ for the most frequent distance between distinct vertices, choosing the largest in a tie as in Fajtlowicz's 1988 convention. Let $\mu(G)$ be the mean over the same unordered pairs, $R(G)=\sum_{uv\in E}(d_ud_v)^{-1/2}$, $\nu(G)$ the maximum number of vertex-disjoint edges, and $D$ the diameter.
\begin{theorem}\label{thm:modecone}
For fixed real $a,b,c\ge0$, the following statements are equivalent:
\begin{enumerate}
\item $m(G)\le a\mu(G)+bR(G)+c\nu(G)$ for every finite connected simple graph on at least two vertices;
\item the inequality holds for every tree having a unique distance mode;
\item $b+c\ge2$.
\end{enumerate}
In particular the sharp coefficient in $m\le\mu+bR$ is $2$. No finite increase of the multiplier of $\mu$ rescues the choice $b=1$, $c=0$.
\end{theorem}
\begin{proof}
Dvořák, Lidický and Škrekovski proved $R\ge D/2$ \cite{dls}. Alternate edges of a diametral path give $\nu\ge D/2$. Since $m\le D$, the third condition implies the first, which implies the second.

For the reverse implication fix an integer $k\ge2$. Start with a path $z_0,\ldots,z_{k-1}$. From each $z_i$ attach a separate path of length $k-i$ to a new vertex $h_i$. All $k$ vertices $h_i$ are at distance $k$ from $z_0$. From $z_0$ grow a separate path of length $L=\lfloor t^{2/3}\rfloor$ ending at $h_*$. Attach $t$ new leaves to each of the $k+1$ vertices $h_0,\ldots,h_{k-1},h_*$. First let $t\to\infty$ at fixed $k$. The diameter is $D=L+k+2$ for all large $t$.

There are $kt^2$ leaf pairs at distance $D$, one leaf at $h_*$ and the other in a local group. Within groups the distance is two, accounting for $(k+1)\binom{t}{2}$ pairs. Between two local groups $i<j$ the distance is $2(k-i)+2$, with $(k-1-i)t^2$ pairs in the bin for $i$. These bin values are distinct and exceed two. The total number of remaining pairs, all touching an original skeleton vertex, is $O_k(tL+L^2)=o_k(t^2)$. Both $k-1$ and $(k+1)/2$ are strictly smaller than $k$. Hence $D$ is the unique mode for all sufficiently large $t$.

Dividing pair distances by $D$, only the remote--local leaf pairs contribute to the limiting mean. There are $(k+1)t+O_k(L)$ vertices, giving
\[
 \frac{\mu}{D}\longrightarrow\frac{2k}{(k+1)^2},\qquad
 \frac{R}{D}\longrightarrow\frac12,\qquad
 \frac{\nu}{D}\longrightarrow\frac12.
\]
For $R$, all but $O(1)$ of the long path's edges contribute $1/2$. The added leaf edges contribute $O_k(\sqrt t)$ and the other core edges $O_k(1)$. For $\nu$, at most $k+1$ matching edges can touch added leaves. Outside those edges a path of length $L$ and $O_k(1)$ other vertices remain. The matching lower bound above supplies the reverse inequality.
The second condition now implies $1\le 2ak/(k+1)^2+(b+c)/2$ for each $k$. Letting $k\to\infty$ proves the third condition.
\end{proof}
The mass being varied here is the number of leaves placed in a group. Pairs between two groups are counted by the product of their masses. Equal depths concentrate many pairs at one distance, while choosing different branching depths distributes the competing pairs among distinct distances. This construction is independent of the stationarity argument for the mean-distance conjecture.

\subsection{Failure even at maximum degree three}
\begin{proposition}\label{thm:cubicmode}
There exist subcubic trees with a unique mode $D-1$ and
\[
 \frac{m}{D}\to1,\qquad \frac{\mu}{D}\to\frac13,
 \qquad \frac{R}{D},\frac{\nu}{D}\to\frac12.
\]
Thus both unit-coefficient conjectures fail on trees of maximum degree three, with normalized deficit tending to $1/6$.
\end{proposition}
\begin{proof}
Join the roots of two complete binary trees of height $h$ by a path of length $L=\lfloor t^{3/2}\rfloor$, with $t=2^h$. Remove one bottom leaf from one tree. The diameter is $D=L+2h$. In the two unpruned trees, inter-tree pairs at distance $D-j$ number at most $(j+1)t^2/2^j$; for $j=0,1$ this is exact. The count follows by writing the two depths as $h-i$ and $h-(j-i)$. After pruning, the counts at $D,D-1$ are $t^2-t$ and $t^2-t/2$ respectively. No pair involving an internal connector vertex reaches $D-1$ when $h\ge2$. For $j\ge2$ the inter-tree bound is at most $3t^2/4$.

For a fixed distance $d$, pairs within one binary tree whose paths branch at a last common ancestor at depth $\ell$ have the following upper bound, summed over both trees:
\[
  \frac{t^2}{2^{\ell+1}}\frac{j_\ell+1}{2^{j_\ell}},
  \qquad j_\ell=2h-2\ell-d\ge0.
\]
There are $2^\ell$ choices of ancestor in each tree; the two opposite child branches have at most $t/2^{\ell+1}$ bottom leaves each. Pairing two levels with a given sum of deficits gives the displayed factor just as above. If $j_\ell<0$ the contribution is zero. At most one $\ell$ has $j_\ell=0$ or 1. At every other depth the last factor is at most $3/4$. Summing the geometric coefficients gives a bound of $7t^2/8$ for these pairs. Ancestor--descendant pairs account for $O(th)$ more. At any one distance, pairs involving an internal connector vertex contribute $O(t+L)$: each vertex in a binary tree has at most one internal connector vertex at that distance, and the connector is a path. For large $h$ the within-tree distances are at most $2h<L$, while all inter-tree distances are at least $L$. Thus all bins other than $D,D-1$ have size at most $7t^2/8+O(th+L)$, strictly below $t^2-t/2$. The unique mode is $D-1$.

The connector contains all but $O(t)$ of the $L+O(t)$ vertices. Its mean pair distance divided by $L$ tends to $1/3$; pairs touching a binary tree constitute $O(t/L)$ of all pairs. Off-connector edges contribute $O(t)$ to $R$ and off-connector vertices can change $\nu$ by $O(t)$, proving the other limits. In any degree-three tree with $D\ge3$, trimming either end of a diametral path gives a path of length $D-1$. Such a shorter path can extend at only one end, or the diameter would be exceeded. At that end there are at most two extensions. Hence the distance bin for $D-1$ is at least as large as the bin for $D$ in every subcubic tree.
\end{proof}

\paragraph{The third distance-mode entry in the 1988 list.}
Item~4 (entry~7 in the later catalogue) asks whether $m\le r+R$, where $r$ is vertex radius. This is true because $r,R\ge D/2$; the diameter bounds already give this implication. The construction in Theorem~\ref{thm:modecone} makes it asymptotically sharp, since the radius of a tree is $\lceil D/2\rceil$. More generally for any $a,b,c,e\ge0$, the inequality $m\le a\mu+bR+c\nu+er$ holds on all connected graphs exactly when $b+c+e\ge2$. Thus the same examples distinguish original items 4, 5 and 6 and locate the precise boundary for all their nonnegative linear combinations.

\section{The exact maximum Wiener index for diameter-five trees}
\label{sec:diamfive}
For a tree \(T\) its Wiener index is \(W(T)=\sum_{\{u,v\}}d_T(u,v)\).
Mukwembi and Vetrík obtained the asymptotically sharp bound
\(9n^2/4-2n^{3/2}+O(n)\) for diameter five
(\href{https://doi.org/10.1017/S0004972713000816}{Theorem 3.3, 2014}).
The exact bound is identified as open in
\href{https://match.pmf.kg.ac.rs/electronic_versions/Match84/n3/match84n3_611-628.pdf}
{MATCH 84 (2020), p.~619}; see also
\href{https://arxiv.org/abs/2605.24855}{Pandey and Ravi, 2026, Section 3}.
We give a formula valid for every \(n\ge6\) and specify every maximizing tree.

\paragraph{Formula.}
Put \(N=n-2\), \(f=n-1\) and let \(t\) be the least positive integer
with \(t(t+1)\ge f\). For \(s\ge2\) define \(\mathcal M(s)\) to consist
of \(\lfloor s/t\rfloor,\lceil s/t\rceil\), each projected to the
integer interval \([1,\lfloor s/2\rfloor]\). For \(m\in\mathcal M(s)\)
put \(q=\lfloor s/m\rfloor,\ r=s-mq\) and
\[
 B(s,m)=(m-r)q^2+r(q+1)^2,\qquad
 H(s)=\min_{m\in\mathcal M(s)}\{fm+B(s,m)\}.
\]
Use all integers \(s\) with \(2\le s\le N/2\) if \(t\le11\); for
\(t\ge12\) use only those with \(N-2s\le19\). Denote this set by \(J\).
Then
\[
 \max_{\substack{|V(T)|=n\\ \operatorname{diam}(T)=5}}W(T)
  =(2n-1)N+\frac{n^2}{4}
    -\min_{s\in J}
         \left\{ \frac{(N-2s)^2}{4}+H(s)+H(N-s)\right\}.
\]
There are at most 64 candidate side splits for \(t\le11\)
(\(n\le133\)) and 10 for \(t\ge12\).
For every minimizing split, join two centers by an edge. At their sides
use masses \(a=s,N-s\); independently choose every minimizing
\(m\in\mathcal M(a)\). Attach \(m-r\) branches of order \(q\) and \(r\)
of order \(q+1\), rooting each at its hub adjacent to the center and
making its remaining vertices leaves. These, up to isomorphism, are
all maximizing trees.

\paragraph{The edge cuts.}
In a diameter-five tree the two centers separate sides of depth two,
each containing a path of length two from its center. Write \(a_1,\ldots,a_m\)
for the orders of the branches on one side (the center is excluded),
and use the same convention on the other side. A branch of order one
temporarily represents a leaf adjacent to a center. An edge whose removal
detaches \(z\) vertices contributes \(z(n-z)\) to \(W(T)\), since precisely
\(z(n-z)\) unordered pairs use it. Denote the two side masses by \(s,N-s\)
and the total number of branches by \(M\). The \(N-M\) leaf edges below
the hubs each contribute \(f\). The hub edges contribute
\(nN-\sum_i a_i^2\) over both sides, and the central edge contributes
\(n^2/4-(N-2s)^2/4\). Thus, apart from the constant in the formula, we
minimize \((N-2s)^2/4+fM+\sum_i a_i^2\).
At fixed side sum and count, replacing two orders \(a<b-1\) by
\(a+1,b-1\) strictly reduces the sum of squares and keeps a branch of
order at least two on that side. Hence optimal branch orders on each
side differ by at most one. If a singleton branch remained, that side
would contain orders one and two. Replacing these by one branch of order
three lowers the penalty by \(f-4=n-5>0\). All optimal orders are therefore
at least two and the least sum of squares is precisely \(B(s,m)\).

\paragraph{Rounding the branch count.}
For fixed \(s\), extend the integer formula for \(fm+B(s,m)\) to positive
real \(m\) using its continuous piecewise-linear expression
\[
      \max_{j\in\mathbb Z_{\ge0}}
              \big\{(2j+1)s+[f-j(j+1)]m\big\}.
\]
Indeed, on \(j\le s/m\le j+1\) the \(j\)-th line dominates, with equality
of neighboring lines at a boundary. Since \(t(t-1)<f\le t(t+1)\), the
function strictly decreases up to \(s/t\) and strictly increases thereafter
when \(f<t(t+1)\).
When \(f=t(t+1)\), its full interval of minimizers is
\([s/(t+1),s/t]\) and it is strictly monotone on either side.
This interval has length \(s/f<1\), as both sides have mass at least two.
Its only possible integer is \(\lfloor s/t\rfloor\). If the interval
contains no integer, the two nearest candidates outside it are
\(\lfloor s/t\rfloor,\lceil s/t\rceil\). Clipping to the feasible range
also handles a minimum outside that range. This proves that \(\mathcal M(s)\)
contains every optimal integer count.

\paragraph{Uniform bound on the number of splits.}
Assume \(t\ge12\), put \(\mu=t+f/t\), and note
\[
       a^2+f-\mu a=(a-t)^2+(t-f/t)(a-t)\ge0
                         \qquad(a\in\mathbb Z_{\ge1}).
\]
Indeed, \(|t-f/t|\le1\), so the squared integer deviation dominates the
linear term. As a competitor take side masses
\(s_i=\lfloor N/2\rfloor,\lceil N/2\rceil\) and choose \(m_i\) nearest
to \(s_i/t\). We have \(N\ge t(t-1)\) and hence
\(m_i\ge t/2-1\), so \(t/(2m_i)\le t/(t-2)<2\).
The balanced branch orders differ from \(t\) by at most two, since
\(|s_i/m_i-t|\le t/(2m_i)\).
Thus every branch contributes at most six in excess of \(\mu a\).
These competitors are feasible, and
\(m_1+m_2\le N/t+1\le t+2\).
The optimal penalty therefore exceeds \(\mu N\) by at most
\(6(t+2)+1/4\le t^2\). For every optimal split this gives
\(|s_i-N/2|\le t\), whence
\(s_i\ge(t^2-3t)/2\).
Each optimal \(m_i\) then satisfies \(|m_i-s_i/t|\le1\)
(the clipping endpoints are inactive) and
\(m_i\ge(t-5)/2\). Consequently
\(|s_i/m_i-t|\le t/m_i\le2t/(t-5)<4\) and all optimal branch
orders belong to \([t-4,t+4]\).
Let \(\Delta=N-2s\ge0\).
If \(\Delta>19\), move a leaf from any branch of order \(y\) on the heavier
side to a branch of order \(x\) on the lighter side. Both sides still have
depth two, and the branch count is unchanged. The penalty changes by
\(1-\Delta+2(x-y+1)\le19-\Delta<0\), a contradiction.
This completes the formula and its equality characterization.

\paragraph{Further consequences and checks.}
At any maximizing tree, no noncentral branch has order $a>n/2$. Otherwise balanced orders force it to be the only branch at its center, since a second branch would have order at least $a-1$, while the other side has mass at least two. Its side mass is then $a$. Move one leaf from it to any branch of order $b\le N-a<a$ on the other side; the penalty changes by $1-(2a-N)+2(b-a+1)<0$. Thus the center on the larger side is a median: each component at that center contains at most $n/2$ vertices. This proves the diameter-five case of Pandey and Ravi's Conjecture~3.10. For a tree, degree distance and Gutman index equal $4W-n(n-1)$ and $4W-(2n-1)(n-1)$ respectively, by summing over edge cuts with degree sums $2z-1,2(n-z)-1$. The same formula and equality graphs therefore determine both maxima.

Distance sums computed by breadth-first search on every diameter-five isomorphism class for $6\le n\le22$ agreed with the formula, including all ties (29,388 trees). A second test generated all 32,500 unrooted trees of orders six through sixteen by their centroid or centroid edge, retained all 2,091 with diameter five, and independently checked the maxima and all equality classes. At $n=12$ the unique maximizer has $W=198$. Its central components have orders 5 and 7: equal component orders need not be optimal.

\section{A self-contained radius bound for the catalogue questions}
\label{sec:radius-certified}
Write $r$ for the vertex radius and give every edge length one. Let $\rho$
be the smallest possible maximum distance from a point on this metric
graph to its vertices. Rounding a center to a nearest vertex gives
$r-1/2\le\rho\le r$. A center can be chosen at a vertex or edge midpoint:
on an edge with coordinate $t\in[0,1]$, the distances to a vertex
are minima of two affine functions with integer intercepts and slopes
$1,-1$. Their ordering is constant on each open half-edge. The maximum
distance is therefore affine there as well, and its minimum occurs at a
half-edge endpoint.

The proof below condenses the degree-cut argument of VMRC~[161,171]
\cite{cat-rad,cat-harm}. We checked the two manuscripts at revision
\texttt{884d261108a87d22}; the corresponding degree-cut and harmonic
bounds are results of those earlier manuscripts.
\begin{theorem}\label{thm:radius-certified}
Every finite connected simple graph satisfies
$M:=\sum_{uv\in E}1/\max(d_u,d_v)\ge\rho$.
If its maximum degree is $\Delta\ge2$, then
\[
 H-r\ge\frac{\Delta-3}{2(\Delta+1)},\qquad
 R-r\ge\sqrt{\Delta}-\tfrac32.
\]
In particular the graphs with $H<r$ and those with $R<r$ are, in both
cases, precisely the even-order paths on at least four vertices.
\end{theorem}

\begin{proof}
\emph{Tree facts.} A tree of diameter $D$ has absolute radius $D/2$:
at the position $i$ on a diameter path, each branch has depth at most
$\min(i,D-i)$, so its vertices lie within $D/2$ of the midpoint.
The ends give the reverse inequality. If a connected $m$-vertex graph
has a vertex incident to $j\ge2$ edges, extend these edges to a spanning
tree. Each of its $j$ branches there contains a different leaf. A longest
path uses at most two leaves, so the tree has diameter at most $m-j+1$.
Consequently the graph satisfies $2r\le m-j+2$.

\emph{Reduction to a core and its leaves.} Suppose $M<\rho$, and minimize
first the cycle rank $|E|-|V|+1$, then $|E|$. Deleting a nonbridge
edge cannot lower $\rho$: distances between retained points can only
increase, and possible center locations are lost. Hence every such deletion
strictly increases $M$ by minimality. Deletion of a leaf on a support
$v$ of degree $d\ge2$ decreases $M$ by
\[
 \frac1d-\frac{b}{d(d-1)}\ge0,
\]
where $b$ counts its other neighbors of degree less than $d$.
The deleted edge contributes $1/d$ and exactly those $b$ edges
increase in weight. This deletion can lower $\rho$ by at most $1/2$.
Indeed, take a vertex or midpoint center $y$ of the remaining graph of
radius $r_0$. If its distance to $v$ is at most $r_0-1/2$, restoration
already costs at most $1/2$. Otherwise that distance equals $r_0$
by half-integrality; move $y$ half a unit toward $v$. This is possible
when the remaining graph has at least two vertices, and covers every
vertex within $r_0+1/2$. Stars and paths are checked directly, including
the singleton and single edge. In a remaining counterexample a support
of degree two has a nonleaf as its other neighbor; here the decrease
of $M$ is $1/2$, contradicting minimality. Two leaves on one support
are also impossible. After removing one, choose a center off the other
pendant interval: projection to the support can only reduce all other
distances, and the remaining leaf is now at distance one, at most
the radius since two vertices remain at distance two. Restoring the
deleted twin costs nothing.

For every $v$, at most one component of $G-v$ contains an edge.
Otherwise split those components into two sides $A,B$, each containing
an edge, and let $G_A$ keep $A\cup\{v\}$ with a new leaf for every
lost neighbor of $v$; define $G_B$ symmetrically. Both graphs preserve
the degree of $v$. Their cycle ranks sum to the original and each
has fewer edges, so both satisfy the desired absolute-radius bound.
Their added leaves together account for exactly one unit of index:
$M(G_A)+M(G_B)=M(G)+1$. Write $r_i=\rho(G_i)$ for $i=A,B$.
Each side has an artificial leaf and a retained neighbor of $v$, so
$r_i\ge1$ and we may project any center on an artificial pendant
interval to $v$. Choose such retained centers $y_i$ and write
$s_i=d(y_i,v)\le r_i-1$ by the distance to an artificial leaf.
On the path from $y_A$ through $v$ to $y_B$, use signed distance
$u\in[-s_A,s_B]$ from $v$. The two sides are covered with radii
$r_A+s_A+u$ and $r_B+s_B-u$, respectively. Their point of intersection,
if on the interval, has value at most $r_A+r_B-1$.
Otherwise the endpoint minimizing the larger function gives radius
$r_A$ or $r_B$, also at most $r_A+r_B-1$. Thus
$\rho(G)\le r_A+r_B-1\le M(G)$, a contradiction.

A tree of diameter at least four violates this cut-vertex property at
the third vertex of a diametral path. A smaller-diameter tree having
at most one leaf per support is a path. Hence our graph contains a
cycle. Choose a maximal induced two-connected subgraph $K$ containing
one. An outside component meeting $K$ at two distinct vertices would
give an ear between them: after deleting any one vertex the union of
the ear with $K$ stays connected. Thus each outside component has a
unique attachment. Removing that attachment shows that the component
contains no edge, since $K$ with one vertex removed is connected and
has an edge. The outside vertices are therefore leaves, with at most
one at each vertex of $K$. Every edge of $K$ is a nonbridge.

\emph{Counting the edges between degree classes.} All degrees in the
next calculation refer to $G$. At a core vertex $v$ let $p_v$ indicate an
attached leaf, $l_v$ count all strictly lower-degree neighbors and put
\[
 g_v=d_v(d_v-1),\qquad
 a_v=l_v/g_v,\qquad b_v=(d_v-l_v)/g_v .
\]
Deleting a core edge $uv$ with $d_u<d_v$ changes $M$ by $a_u-b_v$:
the remaining edges to lower-degree neighbors gain $1/g_u$ or $1/g_v$,
and the removed edge had weight $1/d_v$. For $d_u=d_v=d$ the change is
$(l_u+l_v-d+1)/(d(d-1))$. Both changes must be strictly positive.
If $l_u=0$, then $u$ has no leaf. A higher-degree neighbor violates the
first formula, and an equal-degree neighbor has $l_v\le d_v-1$,
violating the second. Thus $l_v\ge1$. All core degrees are at least
three, as a degree-two vertex has two nonleaf core neighbors. A core
vertex of smallest full degree must support a leaf. Removing all leaves
therefore gives a smaller graph $K$ and
\[
 0\le\delta:=M(G)-M(K)<\rho(G)-\rho(K)\le1 .
\]
The last bound uses the extension of any core center to all attached
leaves. If $\Delta=3$, then every core vertex has one lower-degree
neighbor, its leaf. Deleting any core edge gives change zero, a
contradiction. Hence $\Delta\ge4$.

For $3\le k<\Delta$, let $S$ be the core vertices of full degree at
most $k$, let $N_k$ count the core edges leaving $S$, and orient each
unequal-degree core edge from lower to higher. Denote by $I,F$ the
oriented edges inside and leaving $S$, and by $e_v$ the number of
equal-degree neighbors. At $v$ the indegree is $l_v-p_v$ and the
outdegree is $d_v-l_v-e_v$. Summing the identity
$(d_v-l_v)a_v=l_vb_v$ over $S$ gives
\[
 \sum_{v\in S}p_v b_v
  =\sum_{uv\in I}(a_u-b_v)+\sum_{uv\in F}a_u+\sum_{v\in S}e_v a_v.
\]
Delete only the leaves attached to $S$. The loss of $M$ is the left
side less $\sum_{uv\in Q}1/g_u$, where $Q$ consists of equal-degree
edges with a deleted leaf at each end. Indeed a strict edge changes its
weight only if its higher-degree end loses a leaf, in which case it gains
$1/g_v$ at that end. Its support loses $1/d_v$ and has $l_v-1$ such
core neighbors, leaving precisely $b_v$. An equal-degree edge gains
only when both ends lose leaves.
Each edge of $Q$ contributes $a_u+a_v\ge2/g_u$ to the final sum in the
display; the first sum is positive term by term. Deleting the remaining
leaves can only lower $M$ further. Thus
\[
 \delta\ge\sum_{uv\in F}a_u\ge \frac{N_k}{k(k-1)}.
\]

Write $m=|V(K)|$. On each edge the difference between the reciprocal
degree sum and twice its $M$-weight is $|1/d_u-1/d_v|$. Summing over
edges counts every vertex once in the reciprocal degree sum.
A core edge from degree $d$ to $e>d$ contributes
$\sum_{k=d}^{e-1}1/[k(k+1)]$. Subtract the leaf terms to obtain
\[
 m-2M=\sum_{k=3}^{\Delta-1}\frac{N_k}{k(k+1)}
          -\sum_{v\in K}\frac{p_v}{d_v}\le(\Delta-4)\delta\le\Delta-4.
\]
Here the removed leaf weights sum to at least $\delta$, as retained
edges only gain weight, and the first sum is at most
$\delta\sum_{k=3}^{\Delta-1}(k-1)/(k+1)\le(\Delta-3)\delta$.
A vertex of full degree $\Delta$ has at least $\Delta-1$ neighbors
in $K$. The spanning-tree construction gives a core tree of diameter
at most $m-\Delta+2$. Restoring the leaves increases its diameter by
at most two. Its midpoint yields $2\rho\le m-\Delta+4\le2M$, the
required contradiction.

\emph{Transfer to $H$ and $R$.} Put $n=|V(G)|$ and let an edge have
degrees $a\le b$. Its harmonic and modified weights are
$h=2/(a+b)$ and $m_e=1/b$. The identity
\[
 1/a+1/b=2h+(b/a-1)(h-m_e)
\]
gives $n\le2H+(\Delta-1)(H-M)$. For the ordinary index put
$q=1/\sqrt{ab}$. The analogous identity
$1/a+1/b=2q+(\sqrt{b/a}-1)(q-m_e)$ gives
$n\le2R+(\sqrt\Delta-1)(R-M)$. Using
$M\ge r-1/2$ and $n-2r\ge\Delta-2$ from the spanning-tree
construction yields the two stated bounds. Both are nonnegative for
$\Delta\ge3$. For $\Delta=2$ the graph is a path or a cycle.
For $n\ge3$,
$H(P_n)=(n-3)/2+4/3$, $R(P_n)=(n-3)/2+\sqrt2$ and
$H(C_n)=R(C_n)=n/2$, whereas both graph families have radius
$\lfloor n/2\rfloor$. The one-vertex and two-vertex graphs satisfy
$H=R=r$, completing the classification.
\end{proof}

\section{Graffiti entries 14--16 and Caporossi's temperature refinement}
\label{sec:catalogue}
Fajtlowicz's \emph{Written on the Wall} lists entries 14 and 15 on
printed page~15 and entry 16 on page~18 \cite{cat-wow}. For entry 14
his wording is ``radius is not more than average distance + Randic Index.''
With $r$ denoting radius, $\mu$ the mean distance of unordered distinct
vertices, and $R$ the ordinary Randi\'c index, the three assertions are
\[
 \text{(14)}\quad r\le\mu+R,\qquad
 \text{(15)}\quad r\le V+R,\qquad
 \text{(16)}\quad r\le T+R.
\]
For a graph of order $n$ with vertex degrees $d_v$, the population variance
and mean temperature are
\[
 \overline d=\frac1n\sum_v d_v,\qquad
 V=\frac1n\sum_v(d_v-\overline d)^2,\qquad
 T=\frac1n\sum_v\frac{d_v}{n-d_v}.
\]
The catalogue imposes connectedness for distance questions. We take
$n\ge2$ whenever $\mu$ occurs. Both (15) and (16) are defined also on
$K_1$, where all their terms are zero. Caporossi and Hansen had already
identified the counterexample families for (15) and (16). Fajtlowicz
records their further question about exhaustion on page~15. Caporossi's
thesis confirms that orders 22 and 26 are included and separately
proposes a quantitative temperature bound \cite[pp.~34--35]{cat-thesis}.
The later Caporossi--Hansen paper prints the path predictions for radius
and temperature in Table~1 under entries 12 and 16 \cite{cat-agx2004}.

\begin{theorem}[Complete counterexample classification]\label{thm:cat-cuts}
For every connected simple graph of order at least two,
$R+\mu-r\ge1$ and $H+\mu-r\ge1$, with equality in either bound
exactly on $K_2$. The complete failure sets for entries 15 and 16 and
for their harmonic counterparts are
\[
\begin{array}{c|cc}
 &r\le V+I&r\le T+I\\\hline
 I=R&P_n,\ n\text{ even},\ n\ge22 & P_n,\ n\text{ even},\ n\ge26\\
 I=H&P_n,\ n\text{ even},\ n\ge10 & P_n,\ n\text{ even},\ n\ge14 .
\end{array}
\]
\end{theorem}
\begin{proof}
Theorem~\ref{thm:radius-certified}, proved here from the radius argument of
VMRC~[161,171] \cite{cat-rad,cat-harm}, gives $I\ge r$ for $I=H,R$
except on even paths $P_n$ of order $n\ge4$.
This is the graph-structural input. Both $V$ and $T$ are nonnegative,
and $\mu\ge1$, with equality only on a complete graph. On $K_n$ one has
$H=R=n/2$ and $r=1$ for $n\ge2$. On an even path of order $n\ge4$,
\[
 c_R:=r-R=\tfrac32-\sqrt2,\qquad c_H:=r-H=\tfrac16,
 \qquad \mu=\frac{n+1}{3}.
\]
Indeed its two terminal edges have weights $1/\sqrt2$ and $2/3$ in
$R$ and $H$, respectively. Its $n-3$ remaining edges have weight $1/2$.
There are $n-j$ pairs at distance $j$ for $1\le j<n$, yielding the
formula for $\mu$. Hence $\mu-c_I>1$ for every exceptional path.
On $K_2$ equality holds and $V=0$, $T=1$.

The two ones and $n-2$ twos in the path degree sequence give
\[
 V(P_n)=\frac{2(n-2)}{n^2},\qquad T(P_n)=\frac2{n-1}.
\]
Both decrease strictly for real $n\ge4$ (with $V'$ zero only at~4).
For the ordinary index the adjacent even orders satisfy
\[
 V(P_{22})=\frac{10}{121}<c_R<\frac9{100}=V(P_{20}),\qquad
 T(P_{26})=\frac2{25}<c_R<\frac2{23}=T(P_{24}).
\]
For example, $14142/10000<\sqrt2<14143/10000$, by squaring, puts
$c_R$ strictly between $857/10000$ and $858/10000$.
The corresponding comparisons for the harmonic index are
\[
 V(P_{10})=\frac4{25}<\frac16<\frac3{16}=V(P_8),\qquad
 T(P_{14})=\frac2{13}<\frac16<\frac2{11}=T(P_{12}).
\]
Monotonicity gives precisely the displayed lists.
\end{proof}

\begin{proposition}[Parity resolution of the path refinement]\label{prop:cat-temp}
For every even $n\ge4$ the path uniquely minimizes both $R-r$ and
$T+R-r$ among connected simple graphs of order $n$. Caporossi's
path bound for $R-r$ in equation~(1.16) and his bound for $T+R-r$
in equations~(1.19)--(1.20) therefore hold at every even order.
They also hold at $n=3$. Each fails at every odd order $n\ge5$.
\end{proposition}
\begin{proof}
For either expression the stated bound is its value at $P_n$.
For even $n\ge4$, the radius theorem gives $R-r\ge-c_R$, with the
unique minimizer $P_n$. The path also uniquely minimizes $T$ at every
order $n\ge3$. To see this, choose a spanning tree. Deleting any edge
strictly lowers the temperature sum, since $f(x)=x/(n-x)$ is strictly
increasing. In the tree the degrees are positive integers summing to
$2n-2$. Strict convexity of $f$ shows that replacing a pair $a,b$ with
$a+1,b-1$ for $b-a\ge2$ strictly decreases the sum. Thus the smallest
possible degree sum of $f$ is attained at two ones and $n-2$ twos;
a connected tree with this multiset is a path. At $n=3$ the only
alternative is $K_3$, which has both $R-r=1/2>\sqrt2-1$ and
$T=2>T(P_3)=1$.

For every odd $n\ge5$, attach a new leaf to the neighbor of an endpoint
of $P_{n-1}$ and call the tree $B_n$. Its two short leaves are
interchangeable. Its diameter equals $n-2$, hence its vertex radius
$(n-1)/2$ equals that of $P_n$. The edge types of $B_n$ are two
$(1,3)$ edges, one $(2,3)$ edge, one $(1,2)$ edge and $n-5$ edges
of type $(2,2)$. Thus
\[
 R(P_n)-R(B_n)=1+\frac1{\sqrt2}-\frac2{\sqrt3}-\frac1{\sqrt6}
 >1+\frac7{10}-\frac76-\frac5{12}=\frac7{60}>\frac1{12}.
\]
The degree counts differ from those of $P_n$ by one additional one,
two fewer twos and one additional three. Consequently
\[
 T(B_n)-T(P_n)=\frac{f(1)-2f(2)+f(3)}n
   =\frac{2}{(n-1)(n-2)(n-3)}\le\frac1{12}.
\]
The decrease in $R$ strictly exceeds the increase in $T$. Both
proposed path bounds therefore fail on $B_n$ at each stated order.
\end{proof}

Equations~(1.19)--(1.20) concern a larger right-hand side than entry~16
at odd orders. The odd-order trees in the last proof satisfy entry~16
and refute the path refinement. In the works cited here, the even-path
counterexamples and their thresholds were already documented. The
classification uses the stronger radius theorem to exhaust all graph
shapes. The odd-order refutations above use only the displayed edge and
degree counts.

Table~1 of \cite{cat-agx2004} also considers
$V+R-r\ge n/2-\lfloor n/2\rfloor$ in row~15. Caporossi had already
noted the even-path refutation in his thesis
\cite[p.~35, Conjecture~7]{cat-thesis}. The same $B_n$ supplies
odd-order counterexamples for every odd $n\ge17$. Its degrees give
$V(B_n)=4/n-4/n^2$, decreasing for $n\ge3$. Its edge counts give
$R(B_n)-r=2/\sqrt3+1/\sqrt6+1/\sqrt2-2<271/1000$,
using $2/\sqrt3<11548/10000$, $1/\sqrt6<4083/10000$, and
$1/\sqrt2<7072/10000$, each verified by squaring. Thus $V(B_n)+R(B_n)-r<64/289+271/1000<1/2$.

\section{The exact correction at every odd order}
\label{sec:odd-fork}
For odd $n\ge5$, let $B_n$ be the tree formed by adjoining one leaf
to the neighbor of an end of $P_{n-1}$. Set
\[
 b_R=\frac2{\sqrt3}+\frac1{\sqrt6}+\frac1{\sqrt2}-2,
 \qquad b_H=\frac1{15},\qquad
 t_n=\frac2{n-1}+\frac2{(n-1)(n-2)(n-3)}.
\]
\begin{theorem}\label{thm:odd-fork}
On connected graphs of each odd order $n\ge5$, for $I=R$ and for $I=H$,
\[
       \min (I-r)=b_I,\qquad
       \min (T+I-r)=t_n+b_I .
\]
In all four problems $B_n$ is the unique extremal graph up to isomorphism.
At even orders $n\ge4$ the common unique minimizer is $P_n$.
At $n=2,3$ it is also $P_n$.
\end{theorem}
The earlier fixed-maximum-degree theorem already identifies the relevant
ordinary Randi\'c brooms \cite[paper 05, Section~on the sharp radius bound]{fork-collected}.
The earlier harmonic theorem gives the numerical lower bound $1/15$ and
exhibits these trees \cite{cat-harm}. The proof here identifies every
odd-order harmonic equality case and gives the exact temperature refinements.

\begin{proof}
Fix odd $n\ge5$. The edge and degree counts in
Proposition~\ref{prop:cat-temp} give $R(B_n)-r=b_R$ and $T(B_n)=t_n$.
For $H$, the four exceptional edges contribute $1,2/5,2/3$ in total
for the two short leaf edges, the first long-arm edge and its last edge,
respectively. The other $n-5$ edges have weight $1/2$, so $H(B_n)-r=b_H$.
In particular $b_R<1/2$ and $b_H<1/10$.
If $\Delta\ge4$, Theorem~\ref{thm:radius-certified} gives $R-r\ge1/2$
and $H-r\ge1/10$. Paths and cycles of odd order also lie strictly
above $b_R,b_H$ by their displayed formulas. For $\Delta=3$ the
inequalities before the last substitution in that theorem give, with
$q=\sqrt3$ for $R$ and $q=3$ for $H$,
\[
 (q+1)(I-r)\ge n-2r-(q-1)/2.
\]
Here $n-2r$ is positive and odd by the spanning-tree bound. If it is
at least three then $I-r\ge1/2$ in both cases.
Thus it suffices to consider $n=2r+1$ and $\Delta=3$.

Fajtlowicz's induced-subgraph theorem \cite[Theorem~2]{fork-ciliate}
provides an induced $r$-ciliate. This consists of a cycle of length
$2a$ with a path of length $r-a$ attached at each cycle vertex,
where the two-vertex cycle denotes one edge. The number of vertices is
$2a(r-a+1)=2r+2(a-1)(r-a)$. Thus either $a=1$ or $a=r$,
as only $2r+1$ vertices are available. The induced subgraph is an
even path or an even cycle on $2r=n-1$ vertices. Write $x$ for the
sole other vertex and $k$ for its degree.

Put $Q_I=n-2I$. On an edge of degrees $a,b$, its contribution to $Q_I$
is the sum $1/a+1/b$ minus twice the edge weight, because the reciprocal
degree incidences count every vertex once. Equal degrees contribute zero.
For types $(1,2),(1,3),(2,3)$ denote the contributions by $A,B,C$.
For the two indices they are
\[
\begin{array}{c|ccc}
 &A&B&C\\\hline
 R&(1-1/\sqrt2)^2&(1-1/\sqrt3)^2&(1/\sqrt2-1/\sqrt3)^2\\
 H&1/6&1/3&1/30 .
\end{array}
\]
Both rows satisfy $A>3C>0$ and $B>A+C$.
In the first row $\sqrt B=\sqrt A+\sqrt C$, and the other comparison
reduces to $2\sqrt2>1+\sqrt3$.
For the fork, $Q_I(B_n)=2B+A+C$.
Since $n=2r+1$, minimizing $I-r=(1-Q_I)/2$ amounts to maximizing $Q_I$.

For the induced cycle, if $k=1$ there is one leaf-to-three edge
and two mixed $(2,3)$ edges. For $k=2,3$ there are no leaves and
at most six mixed edges: at most $2k$ lie on the cycle, while $x$
adds at most two if $k=2$ and none if $k=3$. Thus
$Q_I\le\max(B+2C,6C)<2B+A+C$.

For the induced path, $k=1$ and attachment to an endpoint gives a path.
Attachment to an internal vertex gives the fork when that vertex is
adjacent to an endpoint. Otherwise the three leaves produce one
$(1,3)$ edge and two $(1,2)$ edges, and the only degree-three vertex
has two $(2,3)$ edges. The value $B+2A+2C$ is strictly smaller than
the fork value. If $k=2$ or $3$, there are at most two leaves, both
original path ends. Let $\beta$ be the number of $(1,3)$ edges.
For $k=2$, at most two internal path vertices acquire degree three.
Their path incidences supply at most $4-\beta$ mixed $(2,3)$ edges,
and the edges at $x$ supply at most two more.
For $k=3$, let $j$ of its neighbors be original endpoints.
Only the $3-j$ chosen internal vertices have degree three on the path.
Their path incidences supply at most $2(3-j)-\beta$ mixed edges, and
exactly $j$ mixed edges meet $x$. In both cases the mixed count is
at most $6-\beta$. Consequently
\[
 Q_I\le (2-\beta)A+\beta B+(6-\beta)C
       \le 2B+4C<2B+A+C .
\]
This proves uniqueness for $I-r$.

It remains to compare temperatures. Let $f(d)=d/(n-d)$ and let $L$ be
the affine function interpolating $f(1),f(2)$. For every integer
$d\ge3$, convexity gives $f(d)-L(d)\ge f(3)-L(3)>0$.
In a tree the affine sum $\sum_v L(d_v)$ is fixed by the degree sum.
A nonpath tree has some $d_v\ge3$, so its smallest possible sum of
temperatures is attained by the degree multiset of $B_n$, with one
three, three ones and $n-4$ twos. For graphs containing a cycle
the average degree is at least two; Jensen's inequality gives
$T\ge f(2)=2/(n-2)>t_n$ for $n\ge5$.
Hence $B_n$ minimizes $T$ among all nonpaths. It beats the odd path for
$T+R-r$ by Proposition~\ref{prop:cat-temp}; for $T+H-r$ the fork pays
at most $1/12$ extra in $T$ and saves $1/3-1/15>1/12$ in $H-r$.
This proves the odd-order statement. At even orders $n\ge4$, an
even path uniquely minimizes $I-r$ by Theorem~\ref{thm:radius-certified}
and uniquely minimizes $T$ by the convexity argument of
Proposition~\ref{prop:cat-temp}. The cases $K_2$, $P_3$ and $K_3$
give the remaining orders by direct evaluation.
\end{proof}
As a numerical check, all nonisomorphic connected graphs on two through
nine vertices were tested. At orders $5,7,9$ there are respectively
$21$, $853$ and $261080$ graphs. Each of the four optimized quantities
had exactly one minimizing isomorphism class, with the form in
Theorem~\ref{thm:odd-fork}. The proof above covers every order.

\section{A sharp positive gap on every connected nonstar}\label{sec:sixth}
\begingroup\setlength{\emergencystretch}{2em}
Let $H$ denote the harmonic index and $\mu$ the mean distance over distinct vertex pairs. The equality $H=\mu$ holds on stars. There is an order-independent separation for all other graphs.

\paragraph{Weighted-distance input for the harmonic-gap results.}
For the next two harmonic-gap sections we use the weighted distance theorem of the earlier manuscript \cite{weighted}: for every probability $p$ on $V(G)$,
\[
 \sum_{u,v}p_up_v\,d(u,v)+2\sum_v p_v^2\le 1+M,\qquad
 M=\sum_{uv\in E}\frac1{\max(d_u,d_v)} .
\]
Taking the uniform law gives $M\ge[(n-1)\mu-(n-2)]/n$. Putting half the mass on each end of a diameter gives $M\ge D/2$. These are the only conclusions from that theorem needed below.
\begin{theorem}\label{thm:hgap-sixth}
For every finite connected simple graph $G$ which is not a star, $H(G)-\mu(G)\ge1/6$. Equality holds exactly on the four-vertex path.
\end{theorem}
\begin{proof}
For diameter at most two a nonstar is cyclic, so $\mu=2-m/\binom n2$ is at most the value on $S_n^+$, a star with one extra edge. If there is a triangle, the sharp harmonic-index plus girth bound of Wu, Tang and Deng gives $H\ge H(S_n^+)$ \cite{hgapWu}. Without a triangle, a leaf would force a dominating neighbor and hence a star. Thus $\delta\ge2$, $n\ge4$, and disjoint neighborhoods on every edge imply $\Delta\le n-\delta$. Since $2ab/(a+b)^2$ is minimized over $a,b\in[\delta,n-\delta]$ at opposite extremes, summing $2/(a+b)=(1/a+1/b)2ab/(a+b)^2$ gives $H\ge2\delta(n-\delta)/n\ge4-8/n>H(S_n^+)$. In both cases $H-\mu\ge1/2-2(n-3)/[n(n-1)(n+1)]>1/6$. We treat diameter $D\ge3$.

Let $d=\Delta(G)$, $q=\binom n2$ and $L=n-d+1$. A geodesic shows $D\le L$: a vertex outside the geodesic has at most three neighbors on it, and a vertex on it has at most two. Let $B$ be the tree obtained from a path on $L$ vertices by joining $d-1$ extra leaves to one of its endpoints. Let $W$ be the sum of $\min(D,\operatorname{dist}_B(u,v))$ over all unordered distinct pairs on $B$. Then
\begin{equation}\label{eq:hgap-cap}
 \mu(G)\le W/q.
\end{equation}
Indeed extend the star at a vertex of degree $d$ to a spanning tree, building it one leaf at a time. When its current order is $t$, the sorted distances $b_i$ of the new vertex to the $t-1$ earlier vertices satisfy $b_i\le i$. Also $b_i\le t-d+1$, by the geodesic degree bound just proved (the partial tree still contains the original $d$ edges). The broom obtained by extending one fixed arm has precisely $\min(i,t-d+1)$ in position $i$. Distances in the original graph are at most these partial-tree distances and at most $D$. The initial star agrees with the broom construction, proving \eqref{eq:hgap-cap}. Counting leaf--leaf, leaf--path and path--path pairs gives
\begin{align*}
 W={}&(d-1)(d-2)+\frac{(d-1)D(2L-D+1)}2+\frac{LD(D+1)}2\\
 &-\frac{D(D+1)(2D+1)}6+\frac{D(L-D)(L-D-1)}2.
\end{align*}
The formula also holds for $D=L$.

Write $M=\sum_{uv}1/\max(d_u,d_v)$. For $1\le a\le b\le d$ we have $2(d+1)/(a+b)-1/a-d/b=(b-a)(da-b)/[ab(a+b)]\ge0$. Summing yields $(d+1)H\ge n+(d-1)M$. Use
$M\ge[(n-1)\mu-(n-2)]/n$ and $M\ge D/2$ from the uniform and two-point specializations of the weighted theorem. Subtracting $\mu+1/6$ from each lower bound for $H$, both right sides decrease in $\mu$. Substitution of $W/q$ gives
\begin{align*}
 12qn(d+1)(H-\mu-1/6)&\ge U,\\
 12q(d+1)(H-\mu-1/6)&\ge V,
\end{align*}
where
\begin{align*}
 U&=12q[n^2-(d-1)(n-2)]-12(2n+d-1)W-2qn(d+1),\\
 V&=12qn+6q(d-1)D-12(d+1)W-2q(d+1).
\end{align*}
Put $a=d-2$, $b=D-3$, $c=n-d-D+1$. These are nonnegative integers with $a+b+c=n-4$. The two expressions $U,V$ are polynomials of degree at most four in these three variables. The exact rational certificate below proves $\max(U,V)>0$ everywhere except $(0,0,0)$, where $V=0$. The exceptional parameters are $n=4,D=3,d=2$, which force $P_4$. Its $H=11/6$ and $\mu=5/3$ give equality.
\end{proof}

\paragraph{Verification of the quartic signs.}
For $a+b+c<80$ direct substitution covers every triple. For $N=a+b+c\ge80$, write $(a,b,c)=N\,u/s$ with $s=\sum_i u_i$. For each $P\in\{U,V\}$ the polynomial $Q_t(u)=s^4P((80+t)u/s)$ is homogeneous of degree four in $u$ and of degree at most four in $t$. The code triangulates the direction simplex with mesh $1/12$; the vertex coordinates used in the code sum to 12. In each triangle choose $P$ so that $Q_t$ is positive throughout the triangle for every integer $t\ge0$. To see how the finite check proves this, let $v_1,\ldots,v_4$ be any four triangle vertices, repetitions allowed. The symmetric multilinear form associated to $Q_t$ takes the value
\[
 \frac{1}{24}\sum_{\varnothing\ne I\subseteq\{1,2,3,4\}}(-1)^{|I|}
 Q_t\!\left(\sum_{i\in I}v_i\right).
\]
For each such choice the code checks a positive constant term and nonnegative remaining forward differences at $t=0$. Newton's formula $f(t)=\sum_{j=0}^4\binom tj\Delta^jf(0)$ proves positivity at every nonnegative integer $t$. Multilinearity then proves positivity on every convex combination of the triangle vertices, including its boundary. A total of 93 triangles use $U$ and 51 use $V$.
\begin{lstlisting}[basicstyle=\ttfamily\footnotesize,breaklines=true]
"""Exact coverage for the two truncated-broom polynomials; Python 3.9+."""
from fractions import Fraction as F
from functools import cache
from itertools import combinations, combinations_with_replacement

# The arguments are d-2, D-3 and n-d-D+1. The common factor 12 clears halves and sixths.
def bounds(a, b, c):
    d, D, n = a + 2, b + 3, a + b + c + 4
    L = n - d + 1
    q2 = n * (n-1)                 # twice the number of vertex pairs
    W6 = (6*(d-1)*(d-2) + 3*(d-1)*D*(2*L-D+1)
          + 3*L*D*(D+1) - D*(D+1)*(2*D+1)
          + 3*D*(L-D)*(L-D-1))    # six times the truncated broom sum
    U = 6*q2*(n*n-(d-1)*(n-2)) - 2*(2*n+d-1)*W6 - q2*n*(d+1)
    V = 6*q2*n + 3*q2*(d-1)*D - 2*(d+1)*W6 - q2*(d+1)
    return U, V

R, mesh = 80, 12
@cache
def radial_values(u, which):
    s = sum(u)
    return tuple(s**4 * bounds(*(F(t*x,s) for x in u))[which]
                 for t in range(R,R+5))

@cache
def positive_mixed_term(points, which):
    values = [0]*5
    for k in range(1,5):
        for subset in combinations(points,k):
            u = tuple(map(sum,zip(*subset)))
            row = radial_values(u,which)
            for i,v in enumerate(row):
                values[i] += (-1)**k * v
    # The constant in the Newton basis is positive; the other coefficients are nonnegative.
    if values[0] <= 0:
        return False
    while len(values) > 1:
        values = [v-u for u,v in zip(values,values[1:])]
        if values[0] < 0:
            return False
    return True

def covers(triangle, which):
    return all(positive_mixed_term(tuple(sorted(points)),which)
               for points in combinations_with_replacement(triangle,4))

def vertex(i,j):
    return (i,j,mesh-i-j)

counts = [0,0]
for i in range(mesh):
    for j in range(mesh-i):
        B,C = vertex(i+1,j),vertex(i,j+1)
        triangles = [(vertex(i,j),B,C)]
        if i+j < mesh-1:
            triangles.append((B,C,vertex(i+1,j+1)))
        for triangle in triangles:
            which = 1 if covers(triangle,1) else 0
            assert covers(triangle,which)
            counts[which] += 1
assert counts == [93,51]
exceptional = []
for N in range(R):
    for a in range(N+1):
        for b in range(N-a+1):
            c = N-a-b
            u,v = bounds(a,b,c)
            assert max(u,v) >= 0
            if max(u,v) == 0:
                exceptional.append((a,b,c))
assert exceptional == [(0,0,0)]
print('Cover (uniform, diameter):', counts, '| Equality:', exceptional)
\end{lstlisting}
\endgroup

\section{The second-smallest harmonic-index gap at large orders}\label{sec:hgap}
For a connected graph let $H=\sum_{uv\in E}2/(d_u+d_v)$, let $\mu$ be its mean distance over distinct vertex pairs, and put $g=H-\mu$. Write $S_n^+$ for the star with one extra edge joining two leaves. Stars have $g=0$. The next extremum distinguishes the harmonic index from the Randi\'c index.

\begin{theorem}\label{thm:hgap}
For $n\ge10^{60}$ the unique minimizer of $H-\mu$ among connected nonstars of order $n$ is $S_n^+$. Its gap is
\[
 \frac12-\frac4{n(n+1)}+\frac2{n(n-1)}.
\]
In particular, the minimum over nonstars converges to $1/2$.
\end{theorem}
\begin{proof}
Set $M=\sum_{uv\in E}1/\max(d_u,d_v)$, $D=\operatorname{diam}G$ and $x=H-M$. We use the uniform and two-point consequences of the weighted distance theorem:
\[
 M\ge\frac{(n-1)\mu-(n-2)}{n},\qquad M\ge D/2.
\]
In $S_n^+$ the new edge contributes $1/2$, the two affected star edges lose $4/[n(n+1)]$ in total, and the total pairwise distance falls by one. Since $g(S_n^+)<1/2$ for $n>3$, every minimizer satisfies $g<1/2$. Thus $x\le g+1+(\mu-2)/n<5/2$. For an edge with $a\le b$ write $h=2/(a+b)$ and $z=h-1/b$. Then
\begin{equation}\label{eq:hgap-irreg}
 1/a+1/b=2h+(b/a-1)z,\qquad n-2H=\sum_{e}(b/a-1)z.
\end{equation}

\emph{A bound on the mean distance.} Choose a maximal set of vertices of degree at least $\sqrt n$ at mutual distance at least three. Its size $c$ is at most $\sqrt n$, by disjointness of their closed neighborhoods. Let $A$ be the union of the radius-three balls about the chosen vertices, and put $b_0=n-|A|$. Each edge with an endpoint of degree at least $\sqrt n$ lies in $A$. Its contribution to the sum in \eqref{eq:hgap-irreg} is at most $1/a+1/b$; summing these terms costs at most $|A|$. Every other edge costs at most $\sqrt n\,z$. Hence $H\ge b_0/2-(5/4)\sqrt n$.

Contract the components of $G[A]$ and give each new vertex its original size as weight. The other vertices have weight one. A spanning tree of the quotient has $j\le b_0+c$ edges. The sum of the internal diameters of the contracted components is at most $7c$: a connected group of $t$ balls is spanned by traversals of at most $t$ balls of internal diameter six and $t-1$ links. A tree path can therefore be lifted at cost at most $7c$ beyond its length.

For a tree with $j$ edges and positive integer weights summing to $n$, its mean distance for two independent weight-proportional samples is at most $j/2-j^3/(64n^2)+1$. To see this, root at a weighted centroid and write $s_e\le n/2$ for the weight below an edge. The expectation is $j/2-2\sum_e(n/2-s_e)^2/n^2$. For $j\ge8$ there are at most two inclusion-minimal subtrees of weight exceeding $n/2-j/8$, since three disjoint such subtrees weigh more than $n$. Edges with such subtrees form at most two chains, each of length at most $j/8+1$ by integrality. Thus at least $j/2$ edges have deficit at least $j/8$. For $j<8$ the term $+1$ suffices. Conditioning the original uniform samples to be distinct increases their mean distance by at most one. It follows that
\[
 \mu\le j/2-j^3/(64n^2)+7c+2.
\]
Since $H<\mu+1/2$ and $j\le b_0+c$, the preceding estimates imply $j^3/(64n^2)\le(35/4)\sqrt n+5/2$. Consequently $j<9n^{5/6}$ and, at the stated orders,
\[
 \frac{H}{n}<\frac{(9/2)n^{5/6}+7\sqrt n+5/2}{n}<10^{-9},\qquad x<1.501.
\]
Since each edge has harmonic weight at least $1/n$, the number of edges is $m<10^{-9}n^2$.

\emph{A single hub.} Put $\alpha=1/100$. For a vertex of degree $d\ge\alpha n$, at most $2m/(\alpha d)$ of its neighbors have degree greater than $\alpha d$. All its remaining incident edges contribute at least $(1-\alpha)/[(1+\alpha)d]$ to $x$. Hence their total contribution exceeds
\[
 \frac{1-\alpha}{1+\alpha}\left(1-\frac{2m}{\alpha d^2}\right)>0.97.
\]
This contribution is assigned to the unique higher-degree end of the edge. Two vertices of degree at least $\alpha n$ would contradict $x<1.501$. There must be one such vertex, since otherwise \eqref{eq:hgap-irreg} gives $n-2H\le\alpha nx$. Call it $v$, of degree $d$. Its incident edges contribute at most $d+1$ to the sum in \eqref{eq:hgap-irreg}, by summing $1/d+1/d_w$ over $w\sim v$. All other edges contribute at most $\alpha nx$. Thus $k=n-1-d\le2H+\alpha nx<0.0151 n$.

Every pair inside the closed neighborhood of $v$ has distance at most two. The fraction of pairs meeting its complement is at most $2k/n$, so $\mu\le2+0.0302D$. Since $D/2\le H<\mu+1/2$, integrality gives $D\le5$ and $H<2.651$. The edges at $v$ contribute at least
\[
 \sum_{w\sim v}\frac{2}{d+d_w}\ge 2-\frac{2}{d^2}\sum_{w\sim v}d_w\ge2-4m/d^2>1.999.
\]
Consequently the remaining edge weights sum to less than $2/3$. If $G-v$ has maximum degree $t\ge2$, its edges at such a vertex alone contribute at least $t/(t+1)\ge2/3$, since each endpoint degree in $G$ is at most $t+1$. Hence the remaining edges form a matching, and there is at most one, since each weighs at least $1/2$. There must be one because $G$ is a nonstar. Both endpoints must be adjacent to $v$: otherwise connectedness gives endpoint degrees one and two, with edge weight $2/3$. All other vertices must be adjacent to $v$ by connectedness. This gives $G=S_n^+$.
\end{proof}

If $T_n$ is formed by subdividing one edge of $S_{n-1}$, direct calculation gives $g(T_n)=2/3-2/n$. The two candidate values cross between $n=11$ and $n=12$. The theorem does not identify the exact order at which $S_n^+$ first minimizes globally.

\paragraph{Diameter at most two, for every order.}
For every $n\ge3$, $S_n^+$ is the unique minimizer among nonstars with diameter at most two. Indeed, such a graph is cyclic, and therefore $\mu=2-m/\binom n2\le\mu(S_n^+)$. If it contains a triangle, its girth equals three, and the sharp harmonic-index plus girth bound of Wu, Tang and Deng gives $H\ge H(S_n^+)$ with equality only for $S_n^+$ \cite{hgapWu}. Otherwise, a leaf would force its neighbor to dominate the entire graph and triangle-freeness would force a star. Thus $\delta\ge2$ and $n\ge4$. The neighbors of the endpoints of every edge are disjoint, so $\Delta\le n-\delta$. The ratio of $2/(a+b)$ to $1/a+1/b$ is $2ab/(a+b)^2$, minimized over $\delta\le a,b\le n-\delta$ at the two opposite extremes. Summation gives $H\ge2\delta(n-\delta)/n\ge4-8/n$. This exceeds $H(S_n^+)=5/2-2/n-4/[n(n+1)]$ for $n\ge4$.


\section{Which graph metrics have gammoids of maximum-diversity sets?}
\label{sec:block-gammoids}
Let $G$ be a finite connected simple graph, $d$ its shortest-path metric, and
$w:V(G)\to\mathbb R$ a vertex weighting. For $S\subseteq V(G)$ define
\[
 P_w(S)=\sum_{x\in S}w(x)+\sum_{\{x,y\}\subseteq S}d(x,y).
\]
We study \emph{all} subsets maximizing $P_w$ at prescribed cardinality $k$.
A rank-$k$ gammoid here has a directed graph with $k$ designated sources
and terminal set $V(G)$. Its bases are the $k$-sets reachable from these
sources by $k$ pairwise vertex-disjoint directed paths. The source
assignment to the chosen terminals is unrestricted.

\begin{theorem}\label{thm:block-gammoid}
For a connected graph on at least two vertices, the following are equivalent.
\begin{enumerate}
\item Every block of $G$ is a clique.
\item For every $w$ and every $1\le k\le |V(G)|$, the maximizing $k$-sets
for $P_w$ are the bases of a gammoid, realizable in an acyclic network.
\item For every $w$, the maximizing pairs for $P_w$ are the bases of a matroid.
\end{enumerate}
More generally, for an arbitrary finite metric space, the last two
properties are each equivalent to its metric satisfying the four-point
condition.
\end{theorem}

The four-point condition requires that the largest of
\[
 d(a,b)+d(c,d),\qquad d(a,c)+d(b,d),\qquad d(a,d)+d(b,c)
\]
occur at least twice for every four distinct terminals. Buneman's theorem
identifies these metrics with restrictions of weighted tree metrics
\cite{blockbuneman}. Howorka identifies the graph metrics in this class
with block graphs \cite{blockhoworka}. For such a graph the tree is explicit.
Take a node for each vertex and one for each block, and join incident nodes
by an edge of length $1/2$. The resulting incidence tree realizes all the
vertex distances, because crossing any clique block costs one.
Hirai and Murota's characterization of quadratic $M$-convex functions
gives the discrete-convex context for the fixed-rank matroid property
\cite{blockhirai}. The following argument constructs the gammoid from
the tree and gives a four-vertex witness when the metric fails the
condition.

\begin{proof}
Fix a tree realizing the metric, a weighting, and a target rank $k$. Root
the tree. Make internal terminals into leaves by attaching a zero-length
edge carrying the same terminal label. Contract zero-length edges between
internal nodes, and treat an edge ending at a singleton leaf separately
if its length is zero. For each edge $e$ pointing to a child, let $C_e$ be
its descendant terminal set, and let $t_e=|S\cap C_e|$. Counting the
terminal pairs separated by each edge gives, for $|S|=k$,
\[
 P_w(S)=\sum_{x\in S}w(x)+\sum_e \ell_e\,t_e(k-t_e).
\]
For each subtree $C$ let $F_C(j)$ be the largest value for $j$ terminals
of its weight terms plus the displayed edge terms strictly below its
root. The external parameter $k$ is fixed in every term. Take
$0\le j\le\min(k,|C|)$ and $F_C(0)=0$. At a terminal $x$ the single
marginal value is $w(x)$. Suppose a child subtree $C$ is attached to
its parent by an edge of length $\ell$. Offering its $j$th terminal to
the parent contributes the marginal value
\[
 b_{C,j}=F_C(j)-F_C(j-1)+\ell(k-2j+1).
\]
By induction the differences of $F_C$ are nonincreasing. Consequently
each child's offered stream is strictly decreasing if it has at least
two entries: its entering edge has positive length. Merge all the child
streams in nonincreasing order, stopping after $\min(k,|A|)$ entries
for the parent $A$. The sum of the first $t$ entries equals $F_A(t)$.
Indeed, any child allocation chooses a prefix of each stream, and one
can take the largest $t$ entries with arbitrary tie choices, respecting
prefixes by strict decrease within each stream. This also proves the
induction assertion about the parent's differences.

Create one slot vertex for each retained entry at every node. A leaf's
single slot is its labeled terminal. Join a parent slot to every child
slot whose offered value equals its own recorded value. Direct the arcs
toward the leaves. Inductively, paths from the first $t$ slots of a node
to $t$ given leaves exist, with all paths disjoint, exactly when those
leaves attain $F_A(t)$. For the forward implication the paths select the
top $t$ offers with multiplicity. In a child they select a prefix since
all earlier offers are strictly larger. Their continuations certify
optimality inside each child. Conversely an optimal terminal set must
both maximize every used child's subtotal and select top offers at the
parent. Match its occupied child slots value by value to the first $t$
parent slots, and concatenate the child linkages. Different children
are disjoint. Taking the first $k$ slots of the root as sources proves
the gammoid assertion, hence also the ordinary matroid assertion.

Conversely, suppose a quartet has a unique largest pairing sum. Label it
so that $s_1=d(a,b)+d(c,d)$ exceeds both $s_2=d(a,c)+d(b,d)$ and
$s_3=d(a,d)+d(b,c)$. On the quartet put
\[
 w(i)=-\frac12\sum_{j\in\{a,b,c,d\}\setminus\{i\}}d(i,j).
\]
The adjusted scores of $ab,cd$ both equal $-(s_2+s_3)/2$. For $ac,bd$
the score is $-(s_1+s_3)/2$, and for $ad,bc$ it is $-(s_1+s_2)/2$.
Give every outside terminal the same sufficiently negative weight.
The only maximizing pairs are then $ab$ and $cd$, which violate basis
exchange. This proves the necessity and completes both characterizations.
\end{proof}

\begin{corollary}\label{cor:block-positroid}
In every terminal tree metric there is a single ordering of the terminals,
independent of $k,w$, in which each fixed-rank optimum family is a
\emph{positroid}: its bases are the positive maximal minors of a real
full-row-rank matrix whose maximal minors, in increasing column order,
are all nonnegative. This property also characterizes block graphs among
connected graph metrics.
\end{corollary}
\begin{proof}
Make the terminals leaves as above and take their depth-first order.
Write $b_i$ for leaf depth and $u_{ij}$ for minus twice the depth of the
lowest common ancestor of $i,j$. Then $d(i,j)=b_i+b_j+u_{ij}$.
Every clade occupies an interval. Between $i<j$ some consecutive leaves
cross children of their lowest common ancestor, and none crosses above it.
Thus $u_{ij}=\max_{i\le l<j}u_{l,l+1}$. For $t>1$ set
$Y_1=0$, $Y_{l+1}-Y_l=t^{u_{l,l+1}}$, and form the matrix with column
$i$ equal to $t^{w(i)+(k-1)b_i}(1,Y_i,\ldots,Y_i^{k-1})^{\mathsf T}$.
Its Vandermonde minor on $S$ is positive, with leading power
$t^{P_w(S)}$ and positive leading coefficient. Fix a maximizing
$k$-set $B$. Left-multiply by the inverse of its submatrix on $B$.
Every entry now converges as $t\to\infty$: by Cramer's rule the entries
are signed ratios of maximal minors with denominator the minor on $B$,
or entries of the identity. The limiting matrix retains its identity
submatrix, and its ordered maximal minors are the nonnegative limiting
ratios, positive exactly at the maximizing sets. Conversely, positroids
are matroids, so the quartet obstruction still applies.
\end{proof}
The distance sums used here are the weighted dissimilarity vectors
studied through the tropical Veronese Grassmannian by Caminata,
Giansiracusa, Moon and Schaffler \cite{blockcgms}. Speyer and Williams
developed the positive tropical Grassmannian framework \cite{blocksw}.

\paragraph{One network across all ranks.}
The network just constructed may depend on $k$. Even with zero weights,
the quantifier matters. On the full vertex set of a finite connected
graph, the maximizing sets at every rank admit one ordered-source
linking-greedoid representation if and only if the graph is complete.
Indeed, if the diameter is at least two, take diameter endpoints $a,b$
and an internal vertex $x$ of a shortest $a,b$-path. The singleton
$\{x\}$ and the pair $\{a,b\}$ are feasible, whereas neither extension
$\{x,a\}$ nor $\{x,b\}$ is a diameter pair. This violates greedoid
augmentation. In a complete graph all subsets at each rank tie, and a
complete bipartite source-to-terminal network works for every prefix.
Ties of vertex weights or pair distances cause no exceptions in these
arguments.

\paragraph{Exact checks.}
For each connected unlabeled graph on two through seven vertices we
computed its blocks by vertex deletion and tested all four-vertex
weight witnesses in the proof. The block graph counts by order were
$1,2,4,9,22,59$ out of $1,2,6,21,112,853$, respectively; the quartet test
agreed with the block decomposition in every graph. On these block graphs
we compared disjoint-path feasibility against directly computed scores
for all nonempty subsets, using three integer weight profiles (zero,
vertex degree, and a fixed periodic weighting of the vertex labels).
All $27\,705$ comparisons agreed. The proof covers arbitrary real
weightings.

\section{Exact clique-free subgraphs of threshold hosts}\label{sec:threshold-host}
Let $K_s$ denote the complete graph on $s$ vertices. For a host graph $G$, let $\operatorname{ex}(G,K_s)$ be the largest number of edges in a spanning $K_s$-free subgraph. A maximum $r$-cut keeps all host edges whose endpoints belong to different parts of a partition into $r$ classes, with empty classes allowed. Its size is written $b_r(G)$. A threshold graph has a split partition $C\mathbin{\dot\cup}I$, with $C$ complete, $I$ independent, and an ordering $v_1,\ldots,v_c$ of $C$ such that every $x\in I$ has
\[
 N(x)=\{v_1,\ldots,v_{d_x}\},\qquad 0\le d_x\le c .
\]
This follows from the standard nested-neighborhood characterization of threshold graphs. Write $t_q(a)$ for the edge count of the balanced complete $q$-partite graph on $a$ vertices, so $t_1(a)=0$.

\begin{theorem}\label{thm:threshold-turan}
For every threshold graph $G$ and every integer $r\ge2$,
\[
 \operatorname{ex}(G,K_{r+1})=b_r(G)
 =\max_{a\in\{0,\ldots,c\}}
 \left\{t_{r-1}(a)+a(c-a)+\sum_{x\in I}\min(a,d_x)\right\}.
\]
An optimal partition has its first $r-1$ classes balanced on the prefix $\{v_1,\ldots,v_a\}$. The last class consists of every remaining vertex. The same formula holds if each edge from a given $x\in I$ has nonnegative weight $w_x$, with the last summand replaced by $\sum_{x\in I}w_x\min(a,d_x)$ and unit weights on the clique edges.
\end{theorem}
\begin{proof}
Fix a $K_{r+1}$-free subgraph $F$. Choose a largest subset $A\subseteq C$ such that $F[A]$ contains no $K_r$, and put $a=|A|$. For each $v\in C$, its $F$-neighborhood inside $C$ induces a $K_r$-free graph: a forbidden $K_r$ would extend by $v$ to $K_{r+1}$. Thus its degree inside $C$ is at most $a$. Summing those degrees over $C\setminus A$ counts every core edge outside $F[A]$ at least once. There are at most $a(c-a)$ such edges, and Turán's theorem bounds the remaining number by $t_{r-1}(a)$. Likewise the selected neighbors of $x\in I$ form a $K_r$-free subset of $C$, of size at most $\min(a,d_x)$.

For each integer $a$, take the balanced partition of the first $a$ vertices into $r-1$ classes and place all other vertices in one common class. Retaining all cross edges gives exactly the displayed objective: the prefix condition ensures that precisely $\min(a,d_x)$ neighbors of $x$ lie in the first $r-1$ classes. This construction is $r$-partite and therefore $K_{r+1}$-free. Nonnegative outer weights preserve the upper bound and the construction's equality.
\end{proof}

If $F_r(a)$ is the objective, its forward difference is
\[
 F_r(a+1)-F_r(a)
  =c-a-1-\left\lfloor\frac a{r-1}\right\rfloor+\sum_{d_x>a}w_x
 \qquad(0\le a<c).
\]
Indeed, adding a vertex to a smallest part of a balanced $(r-1)$-partition adds $a-\lfloor a/(r-1)\rfloor$ edges. The difference decreases strictly. Therefore the maximum occurs at one index or at two consecutive indices when a difference vanishes. Histograms of the values $d_x$ and their suffix sums give the value and a witnessing partition in $O(c+|I|)$ arithmetic operations once the ordered representation is supplied. This proves existence of an optimal partition of the indicated form and does not assert that every extremal subgraph has it. Polynomial-time Max Cut on threshold graphs already follows from the cograph algorithm of Bodlaender and Jansen \cite{bjcut}; Dross et al.\ give a linear-time algorithm on an ordered threshold graph for the equivalent complementary completion problem \cite[Theorem~4.6]{drosscompletion}. The result here simultaneously optimizes over all clique-free subgraphs and all $r$-partitions.

\subsection{The obstruction on general split graphs}
The construction in the following proposition is the split-graph reduction of Bodlaender and Jansen \cite{bjcut}. Its Max Cut identity is also reproduced by Faria et al.\ \cite[Theorem~5]{fariacut}. The same construction gives exact offsets for three other triangle parameters. Write $\tau(J)$ for the minimum number of edges meeting all triangles and $\nu(J)$ for the maximum number of edge-disjoint triangles.
\begin{proposition}\label{prop:split-offset}
Let $J$ have vertex set $C$ and $s$ missing pairs. Complete $C$ to a clique and, for each missing pair $uv$, adjoin one vertex adjacent precisely to $u,v$, with all new vertices independent. Denote the resulting split graph by $G$. Then
\begin{align*}
 b_2(G)&=2s+b_2(J),&
 \operatorname{ex}(G,K_3)&=2s+\operatorname{ex}(J,K_3),\\
 \tau(G)&=s+\tau(J),&
 \nu(G)&=s+\nu(J).
\end{align*}
\end{proposition}
\begin{proof}
The new edges partition into $s$ edge-disjoint triangles, each consisting of a completed pair and its two new incident edges. A cut keeps at most two edges of each such triangle. For any prescribed two-coloring of $C$, the new vertices can be colored individually to keep exactly two. This proves the first identity.

A triangle-free subgraph also keeps at most two edges per new triangle, and its original edges must themselves be triangle-free. Conversely retain a maximum triangle-free subgraph of $J$, discard all completed pairs and keep every edge at a new vertex. No new vertex lies on a triangle of the resulting subgraph. This proves the second identity. Since $|E(G)|=|E(J)|+3s$ and $\tau=|E|-\operatorname{ex}(\,\cdot\,,K_3)$, the third identity follows. Any packed triangle outside $J$ uses at least one completed pair, so at most $s$ packed triangles can lie outside $J$. All $s$ new triangles can accompany any packing of $J$, proving the last equality.
\end{proof}
Taking $J$ to be the disjoint union of $t$ pentagons gives a connected split host $G$ with $\operatorname{ex}(G,K_3)-b_2(G)=t$. Thus the equality of the two extremal values in Theorem~\ref{thm:threshold-turan} fails by arbitrarily large additive amounts on split hosts, even when every vertex of the independent part has degree two.

\begin{thebibliography}{99}
\bibitem{mode-faj} S. Fajtlowicz, On conjectures of Graffiti, \emph{Discrete Mathematics} \textbf{72} (1988), 113--118, list of open conjectures, item~5. \href{https://doi.org/10.1016/0012-365X(88)90199-9}{Publisher record}.

\bibitem{cat-wow} S. Fajtlowicz, \emph{Written on the Wall}, electronic catalogue, July 2004 copy, pp.~14, 15, 18, entries 8--9 and 14--16. \href{https://independencenumber.wordpress.com/wp-content/uploads/2012/08/wow-july2004.pdf}{Catalogue PDF}.
\bibitem{cat-thesis} G. Caporossi, \emph{D\'ecouverte par ordinateur en th\'eorie des graphes}, Ph.D. thesis, \'Ecole Polytechnique de Montr\'eal, Universit\'e de Montr\'eal, August 2000, pp.~33--35, equations~(1.14)--(1.20). \href{https://www.collectionscanada.gc.ca/obj/s4/f2/dsk3/ftp05/NQ65536.pdf}{National Library of Canada PDF}.
\bibitem{cat-rad} Verifiable Mathematical Reasoning Corpus, \emph{A proof of the Caporossi--Hansen conjecture on the Randi\'c index and graph radius}, manuscript [161], November 2026. \href{https://github.com/QEEDDD/Verifiable-Mathematical-Reasoning-Corpus/blob/884d261108a87d2229e5bcc273e5f4a607c771d2/\%5B161\%5D\%20Randic_Index_and_Vertex_Radius.tex}{Frozen source manuscript}.
\bibitem{cat-harm} Verifiable Mathematical Reasoning Corpus, \emph{Harmonic index and radius: a strengthening of the Caporossi--Hansen inequality}, manuscript [171], November 2026. \href{https://github.com/QEEDDD/Verifiable-Mathematical-Reasoning-Corpus/blob/884d261108a87d2229e5bcc273e5f4a607c771d2/\%5B171\%5D\%20Harmonic_Radius_Master_Theorem.tex}{Frozen source manuscript}.
\bibitem{cat-agx2004} G. Caporossi and P. Hansen, Variable neighborhood search for extremal graphs. 5. Three ways to automate finding conjectures, \emph{Discrete Mathematics} \textbf{276} (2004), 81--94, Table~1. \href{https://doi.org/10.1016/S0012-365X(03)00311-X}{Publisher record}.

\bibitem{weighted} Verifiable Mathematical Reasoning Corpus, \emph{The Randi\'c index is at least the average distance: a proof of Fajtlowicz's conjecture}, manuscript, 20 November 2026, Theorem~1.1. See also Shaitetsu, \emph{A proof audit of Fajtlowicz's average-distance conjecture}, 21 November 2026, Secs.~2--3.
\bibitem{dls} Z. Dvořák, B. Lidický and R. Škrekovski, Randi\'c index and the diameter of a graph, \emph{European J. Combin.} \textbf{32} (2011), 434--442. \doi{10.1016/j.ejc.2010.12.002}.
\bibitem{mukwet} S. Mukwembi and T. Vetrík, Wiener index of trees of given order and diameter at most 6, \emph{Bull. Aust. Math. Soc.} \textbf{89} (2014), 379--396. \doi{10.1017/S0004972713000816}.
\bibitem{pandey} D. Pandey and P. S. Ravi, On a conjecture of DeLaViña and Waller, arXiv:2605.24855 (2026). \url{https://arxiv.org/abs/2605.24855}.

\bibitem{fork-ciliate} S. Fajtlowicz, A characterization of radius-critical graphs, \emph{Journal of Graph Theory} \textbf{12} (1988), 529--532. \href{https://doi.org/10.1002/jgt.3190120409}{doi:10.1002/jgt.3190120409}.
\bibitem{fork-collected} Keio and Shaitetsu, \emph{Collected research papers: the Randi\'c index, graph distances, and rational Lambert series}, VMRC [189], paper 05. \href{https://github.com/QEEDDD/Verifiable-Mathematical-Reasoning-Corpus/blob/884d261108a87d2229e5bcc273e5f4a607c771d2/\%5B189\%5D\%20Collected_Research_Papers.tex}{Frozen source manuscript}.

\bibitem{hgapWu} R. Wu, Z. Tang and H. Deng, On the harmonic index and the girth of a graph, \emph{Utilitas Math.} \textbf{91} (2013), 65--69. \url{https://utilitasmathematica.com/index.php/Index/article/view/970}.
\bibitem{hgapShi} L. Shi, Chemical indices, mean distance, and radius, \emph{MATCH Commun. Math. Comput. Chem.} \textbf{75} (2016), 57--70. \url{https://match.pmf.kg.ac.rs/electronic_versions/Match75/n1/match75n1_57-70.pdf}.

\bibitem{blockbuneman}
P. Buneman, A note on the metric properties of trees,
\emph{J. Combin. Theory Ser. B} \textbf{17} (1974), 48--50.
\href{https://doi.org/10.1016/0095-8956(74)90047-1}{doi:10.1016/0095-8956(74)90047-1}.
\bibitem{blockhoworka}
E. Howorka, On metric properties of certain clique graphs,
\emph{J. Combin. Theory Ser. B} \textbf{27} (1979), 67--74.
\href{https://doi.org/10.1016/0095-8956(79)90069-8}{doi:10.1016/0095-8956(79)90069-8}.
\bibitem{blockhirai}
H. Hirai and K. Murota, $M$-convex functions and tree metrics,
\emph{Japan J. Indust. Appl. Math.} \textbf{21} (2004), 391--403.
\href{https://doi.org/10.1007/BF03167590}{doi:10.1007/BF03167590}.
\bibitem{blockcgms}
A. Caminata, N. Giansiracusa, H.-B. Moon and L. Schaffler,
Point configurations, phylogenetic trees, and dissimilarity vectors,
\emph{Proc. Natl. Acad. Sci. USA} \textbf{118} (2021), e2021244118.
\href{https://doi.org/10.1073/pnas.2021244118}{doi:10.1073/pnas.2021244118}.
\bibitem{blocksw}
D. Speyer and L. Williams, The tropical totally positive Grassmannian,
\emph{J. Algebraic Combin.} \textbf{22} (2005), 189--210.
\href{https://arxiv.org/abs/math/0312297}{arXiv:math/0312297}.

\bibitem{bjcut} H.~L.~Bodlaender and K.~Jansen, On the complexity of the maximum cut problem, \emph{Nordic J. Comput.} \textbf{7} (2000), 14--31.
\bibitem{drosscompletion} F.~Dross, C.~Hilaire, I.~Koch, V.~Leoni, N.~Pardal, M.~I.~Lopez Pujato and V.~F.~dos Santos, On the proper interval completion problem within some chordal subclasses, \emph{Discrete Math.} \textbf{348} (2025), 114274. \url{https://doi.org/10.1016/j.disc.2024.114274}
\bibitem{fariacut} L.~Faria, S.~Klein, I.~Sau and R.~Sucupira, Improved kernels for Signed Max Cut parameterized above lower bound on $(r,\ell)$-graphs, \emph{Discrete Math. Theor. Comput. Sci.} \textbf{19} (2017), no.~1. \url{https://dmtcs.episciences.org/en/articles/3682}

\end{thebibliography}
\end{document}